% =====================================================================
%  第 5 章 几何量子化 —— ch05b（§5.3–§5.4）
%  页段：PDF 273–297（原书 p.272–296）
% =====================================================================

% [原文 p.272]
% [块界：本句起于 PDF p.272 中段，已按前页图补全；上一块 ch05a 已在例 5.2.4 中译出「水平集是拉格朗日子流形」，本块从「我们有」续起，不重复]
我们有
\[
\frac{1}{2\pi}\oint_C \lambda=\frac{1}{2\pi}\oint_C p\,\rd x=\frac{E}{\omega},
\qquad \omega=\sqrt{\frac{k}{m}},\quad m(C)=2.
\]
凯勒–马斯洛夫量子化条件为
\[
\frac{E}{\hbar\omega}-\frac{1}{2}\in\Z.
\]
这与众所周知的能谱
\[
E=\hbar\omega\left(n+\frac{1}{2}\right),\qquad n=0,1,2,\cdots
\]
相容。

\takeaway{§5.2 把 WKB 展开重写成几何形式：在 $\hbar$-渐近解 $\psi=e^{\frac{i}{\hbar}S}a$ 中，相函数 $S$ 满足哈密顿–雅可比方程、其等值面即拉格朗日子流形；半经典波能否在拉格朗日子流形上整体定义，取决于回路上 $\frac{1}{2\pi\hbar}\oint\lambda$ 与马斯洛夫修正能否凑成整数，这就是凯勒–马斯洛夫（EBK）量子化条件，谐振子一例给出精确能谱。}

\section{海森堡群}

\subsection{海森堡李代数}

设 $(V,\omega)$ 是 $2n$ 维辛向量空间。考虑 $2n+1$ 维向量空间
\[
h_V:=V\oplus\R K,
\]
其中我们用 $K$ 表示新增维度的基向量。我们可以用下式在 $h_V$ 上定义括号：
\[
[u,v]=\omega(u,v)K\quad (u,v\in V),\qquad [-,K]=0,
\]
即 $K$ 是中心元。
容易验证 $(h_V,[-,-])$ 构成李代数，称为\textbf{海森堡李代数}（Heisenberg Lie algebra）。具体地，设 $\{e_1,\cdots,e_n,f_1,\cdots,f_n\}$ 是 $(V,\omega)$ 的辛基，则
\[
\begin{gathered}
[e_i,f_j]=\delta_{ij}K,\\
[e_i,e_j]=[f_i,f_j]=0,\\
[e_i,K]=[f_i,K]=0
\end{gathered}
\]
给出海森堡李代数中显式的括号关系。

当 $V=\R^{2n}$ 是标准辛空间时，记
\[
h_n:=h_{\R^{2n}}.
\]
由于任意 $2n$ 维辛向量空间 $(V,\omega)$ 都辛同构于标准空间 $\R^{2n}$，我们有李代数同构
\[
h_V\simeq h_n.
\]
设 $\varphi\in\operatorname{Sp}(V)$ 是辛群中的元素。由于 $\varphi$ 保持辛配对，它诱导李代数同构
\[
\varphi:\ h_V\longrightarrow h_V,\qquad
u+\lambda K\longmapsto \varphi(u)+\lambda K.
\]
由此得到群同态
\[
\operatorname{Sp}(V)\longrightarrow \Aut(h_V).
\]

% [原文 p.273]
\subsection{海森堡群}

海森堡群 $H_V$ 是与李代数 $h_V$ 相联系的单连通李群。作为集合，$H_V$ 可通过指数映射与 $h_V$ 等同：
\[
H_V=\left\{e^{u+\lambda K}\right\}_{u+\lambda K\in h_V}.
\]
乘法法则可由贝克–坎贝尔–豪斯多夫公式得到：
\[
e^Ae^B=e^{A+B+\frac{1}{2}[A,B]+\text{更高阶括号}}.
\]
由于 $[h_V,h_V]=\R K$ 是中心，这里所有更高阶括号都消失。因此
\[
e^{u_1+\lambda_1 K}\cdot e^{u_2+\lambda_2 K}
=e^{u_1+u_2+\left(\lambda_1+\lambda_2+\frac{1}{2}\omega(u_1,u_2)\right)K},
\qquad u_1,u_2\in V.
\]
当 $V=\R^{2n}$ 是标准辛空间时，记
\[
H_n:=H_{\R^{2n}}.
\]
对一般的 $2n$ 维辛向量空间 $V$，选定一组辛基便给出群同构
\[
H_V\simeq H_n.
\]
此外还有群同态
\[
\operatorname{Sp}(V)\longrightarrow\Aut(H_V).
\]
给定 $\varphi\in\operatorname{Sp}(V)$，诱导的群同构为
\[
H_V\longrightarrow H_V,\qquad
e^{u+\lambda K}\longmapsto e^{\varphi(u)+\lambda K}.
\]
这与海森堡李代数上相应的变换相容。

\subsection{薛定谔表示}

设 $V=\R^{2n}$，其标准辛基为 $\{e_1,\cdots,e_n,f_1,\cdots,f_n\}$。令
\[
\mathcal{S}(\R^n)=\left\{f(\mathbf x)\mid f \text{ 是 } \R^n \text{ 上的 } \C \text{ 值施瓦茨函数}\right\}.
\]
$\mathcal{S}(\R^n)$ 通过下式给出海森堡李代数 $h_n$ 的一个表示：
\[
e_i\longmapsto \frac{i}{\hbar}\hat x^i=\frac{i}{\hbar}x^i,\qquad
f_i\longmapsto -\frac{i}{\hbar}\hat p_i=-\frac{\pa}{\pa x^i},\qquad
K\longmapsto \frac{i}{\hbar}.
\]
这可以指数化，成为海森堡群 $H_n$ 的酉表示。事实上，我们把 $h_n$ 中的向量写成
\[
\mathbf r\cdot\mathbf e+\mathbf s\cdot\mathbf f+\lambda K.
\]
这里 $\mathbf r=(r_1,\cdots,r_n)$、$\mathbf s=(s_1,\cdots,s_n)\in\R^n$，而

% [原文 p.274]
\[
\mathbf e=(e_1,\cdots,e_n),\qquad \mathbf f=(f_1,\cdots,f_n)
\]
是上述辛基。

给定函数 $\psi(\mathbf x)\in\mathcal{S}(\R^n)$，我们希望定义变换后的函数
\[
e^{\mathbf r\cdot\mathbf e+\mathbf s\cdot\mathbf f+\lambda K}\psi.
\]
为说明它如何定义，我们先利用群乘法法则改写
\[
e^{\mathbf r\cdot\mathbf e+\mathbf s\cdot\mathbf f+\lambda K}
=e^{\mathbf r\cdot\mathbf e+\left(\lambda-\frac{1}{2}\mathbf r\cdot\mathbf s\right)K}
e^{\mathbf s\cdot\mathbf f}.
\]
由相应的李代数作用，以及 $\pa_{x^i}$ 生成平移这一事实，
\[
\left(e^{\mathbf s\cdot\mathbf f}\psi\right)(\mathbf x)
=\left(e^{-\mathbf s\cdot\nabla}\psi\right)(\mathbf x)
=\psi(\mathbf x-\mathbf s),
\]
\[
e^{\mathbf r\cdot\mathbf e+\left(\lambda-\frac{1}{2}\mathbf r\cdot\mathbf s\right)K}\psi(\mathbf x)
=e^{\frac{i}{\hbar}\left(\mathbf r\cdot\mathbf x-\frac{1}{2}\mathbf r\cdot\mathbf s+\lambda\right)}\psi(\mathbf x).
\]
因此预期的变换规律为
\[
\left(e^{\mathbf r\cdot\mathbf e+\mathbf s\cdot\mathbf f+\lambda K}\psi\right)(\mathbf x)
=e^{\frac{i}{\hbar}\left(\mathbf r\cdot\mathbf x-\frac{1}{2}\mathbf r\cdot\mathbf s+\lambda\right)}
\psi(\mathbf x-\mathbf s).
\]

\begin{definition}{5.3.1}
  给定 $(\mathbf r,\mathbf s,\lambda)\in\R^{2n+1}$，我们定义
  \[
  T_\hbar(\mathbf r,\mathbf s,\lambda):\ \mathcal{S}(\R^n)\longrightarrow\mathcal{S}(\R^n)
  \]
  为
  \[
  (T_\hbar(\mathbf r,\mathbf s,\lambda)\psi)(\mathbf x)
  :=e^{\frac{i}{\hbar}\left(\mathbf r\cdot\mathbf x-\frac{1}{2}\mathbf r\cdot\mathbf s+\lambda\right)}
  \psi(\mathbf x-\mathbf s).
  \]
\end{definition}

\begin{proposition}{5.3.2}
  如下复合法则成立：
  \[
  T_\hbar(\mathbf r_1,\mathbf s_1,\lambda_1)\,T_\hbar(\mathbf r_2,\mathbf s_2,\lambda_2)
  =T_\hbar\left(\mathbf r_1+\mathbf r_2,\ \mathbf s_1+\mathbf s_2,\ 
  \lambda_1+\lambda_2+\frac{1}{2}(\mathbf r_1\cdot\mathbf s_2-\mathbf r_2\cdot\mathbf s_1)\right).
  \]
\end{proposition}

这可由上面的形式计算得到，也可直接验证。于是 $T_\hbar$ 给出海森堡群 $H_n$ 的表示
\[
e^{\mathbf r\cdot\mathbf e+\mathbf s\cdot\mathbf f+\lambda K}\psi
:=T_\hbar(\mathbf r,\mathbf s,\lambda)\psi.
\]
显然 $T_\hbar$ 保持厄米内积。因此我们可以把 $T_\hbar$ 延拓为
\[
T_\hbar(\mathbf r,\mathbf s,\lambda):\ L^2(\R^n)\longrightarrow L^2(\R^n).
\]
$T_\hbar$ 给出 $H_n$ 在 $L^2(\R^n)$ 上的酉表示。于是有群同态
\[
T_\hbar:\ H_n\longrightarrow U(L^2(\R^n)),
\]
其中 $U(L^2(\R^n))$ 是 $L^2(\R^n)$ 的酉群。这称为 $H_n$ 的以 $\hbar$ 为参数的\textbf{薛定谔表示}（Schrödinger representation）。该参数体现中心元 $K$ 的取值
\[
K\longmapsto \frac{i}{\hbar}.
\]

% [原文 p.275]
酉表示 $T_\hbar$ 通过下式给出量子算子的变换：
\[
\mathrm{Ad}_{T_\hbar}:\ \wh\Theta\longmapsto T_\hbar\wh\Theta T_\hbar^{-1}.
\]
例如
\[
\mathrm{Ad}_{T_\hbar(\mathbf r,\mathbf s,\lambda)}(\wh x^i)=\wh x^i-s_i,
\qquad
\mathrm{Ad}_{T_\hbar(\mathbf r,\mathbf s,\lambda)}(\wh p_i)=\wh p_i-r_i.
\]
可以证明，海森堡群 $H_n$ 在 $L^2(\R^n)$ 上的酉表示 $T_\hbar$ 是不可约的，即 $L^2(\R^n)$ 中在所有 $H_n$ 作用下不变的闭子空间只有 $\{0\}$ 与 $L^2(\R^n)$ 本身。此外，著名的斯通–冯·诺依曼定理断言：这本质上是 $H_n$ 在希尔伯特空间上唯一的不可约酉表示。

\begin{theorem}{5.3.3}
  （斯通–冯·诺依曼）设 $\rho$ 是 $H_n$ 在希尔伯特空间 $\mathbb{H}$ 上任意的不可约酉表示，其中心作用为 $\rho(e^K)=e^{\frac{i}{\hbar}}$。则存在酉算子 $U:\mathbb{H}\to L^2(\R^n)$ 使得
  \[
  U\rho U^{-1}=T_\hbar.
  \]
  这样的 $U$ 在相差相位因子 $\xi$（$|\xi|=1$）的意义下唯一。
\end{theorem}

证明见定理 4.4.1。

\subsection{巴格曼–福克表示}

还有另一种用全纯函数得到海森堡群不可约表示的方式。它与福克空间的构造密切相关。

首先，我们在 $\C^n$ 上的两个函数之间定义厄米内积
\[
\langle f|g\rangle_\hbar
:=\frac{1}{(\pi\hbar)^n}\int_{\C^n}\rd^{2n}z\,
\overline{f(z)}g(z)\,e^{-\abs{z}^2/\hbar}.
\]
这里 $f,g$ 不必全纯，不过下文讨论 Segal–Bargmann 空间时我们会限制到全纯函数。归一化常数选取为使
\[
\langle 1|1\rangle_\hbar=1.
\]
相应的 $L^2$-范数记为
\[
\norm{f}_\hbar^2=\langle f|f\rangle_\hbar
=\frac{1}{(\pi\hbar)^n}\int_{\C^n}\rd^{2n}z\,\abs{f(z)}^2e^{-\abs{z}^2/\hbar}.
\]

\begin{definition}{5.3.4}
  \textbf{Segal–Bargmann 空间}（或称巴格曼–福克空间）定义为
  \[
  \mathcal{F}_n:=\left\{f(z)\ \text{在}\ \C^n\ \text{中全纯，且}\ \norm{f}_\hbar^2<\infty\right\}.
  \]
\end{definition}

\begin{proposition}{5.3.5}
  $\mathcal{F}_n$ 是希尔伯特空间。
\end{proposition}

% [原文 p.276]
\begin{proof}
  我们需要证明 $\mathcal{F}_n$ 关于范数 $\norm{\cdot}_\hbar$ 完备。设 $P_R(\mathbf z_0)$ 表示以 $\mathbf z_0$ 为心、半径为 $R$ 的多圆盘：
  \[
  P_R(\mathbf z_0)=\left\{\mathbf z\in\C^n\mid \abs{z^i-z_0^i}<R,\ i=1,\cdots,n\right\}.
  \]
  对任意全纯函数 $f(\mathbf z)$，我们有平均值性质
  \[
  f(\mathbf z_0)=\frac{1}{(\pi R^2)^n}\int_{P_R(\mathbf z_0)}\rd^{2n}\mathbf z\, f(\mathbf z).
  \]
  用柯西–施瓦茨不等式，
  \[
  \begin{aligned}
  \abs{f(\mathbf z_0)}
  &\le\frac{1}{(\pi R^2)^n}\int_{P_R(\mathbf z_0)}\rd^{2n}\mathbf z\,
  \abs{f(\mathbf z)}e^{-\abs{\mathbf z}^2/2\hbar}e^{\abs{\mathbf z}^2/2\hbar}\\
  &\le\frac{1}{(\pi R^2)^n}
  \left(\int_{P_R(\mathbf z_0)}\rd^{2n}\mathbf z\, e^{\abs{\mathbf z}^2/\hbar}\right)^{\frac{1}{2}}
  \left(\int_{P_R(\mathbf z_0)}\rd^{2n}\mathbf z\,\abs{f(\mathbf z)}^2e^{-\abs{\mathbf z}^2/\hbar}\right)^{\frac{1}{2}}\\
  &\le\frac{(\pi\hbar)^n}{(\pi R^2)^n}
  \left(\int_{P_R(\mathbf z_0)}\rd^{2n}\mathbf z\, e^{\abs{\mathbf z}^2/\hbar}\right)^{\frac{1}{2}}
  \norm{f}_\hbar.
  \end{aligned}
  \]
  \tutornote{末行系数 $(\pi\hbar)^n$ 按量纲似应为 $(\pi\hbar)^{n/2}$；此处照原文录入。}
  这表明：对 $\mathbf z_0$ 的任一开邻域 $V$，可以选取依赖于 $V$ 的常数 $C_V$，使得
  \[
  \abs{f(\mathbf z)}\le C_V\norm{f}_\hbar,\qquad \forall\mathbf z\in V
  \]
  对任意 $f\in\mathcal{F}_n$ 成立。
  现设 $\{f_n\}$ 是 $\mathcal{F}_n$ 中的柯西列。上述估计表明，$\{f_n\}$ 局部一致收敛到某个极限函数 $f$，而 $f$ 必定也全纯。于是 $f\in\mathcal{F}_n$。因此 $\mathcal{F}_n$ 完备，即它是希尔伯特空间。
\end{proof}

\begin{proposition}{5.3.6}
  设 $f,g$ 是关于 $\mathbf z=(z^1,\cdots,z^n)$ 的多项式函数。则
  \[
  \langle z^if|g\rangle_\hbar=\langle f|\hbar\pa_{z^i}g\rangle_\hbar.
  \]
\end{proposition}

\begin{proof}
  \[
  \begin{aligned}
  \langle z^if|g\rangle_\hbar
  &=\frac{1}{(\pi\hbar)^n}\int_{\C^n}\rd^{2n}\mathbf z\,
  \overline{z^if(\mathbf z)}g(\mathbf z)e^{-\abs{\mathbf z}^2/\hbar}\\
  &=\frac{1}{(\pi\hbar)^n}\int_{\C^n}\rd^{2n}\mathbf z\,
  \overline{f(\mathbf z)}g(\mathbf z)\left(-\hbar\frac{\pa}{\pa z^i}
  e^{-\abs{\mathbf z}^2/\hbar}\right)\\
  &=\frac{1}{(\pi\hbar)^n}\int_{\C^n}\rd^{2n}\mathbf z\,
  \overline{f(\mathbf z)}\left(\hbar\pa_{z^i}g(\mathbf z)\right)
  e^{-\abs{\mathbf z}^2/\hbar}\\
  &=\langle f|\hbar\pa_{z^i}g\rangle_\hbar.
  \end{aligned}
  \]
\end{proof}

我们定义升算子与降算子 $a_i^\dagger$ 和 $a_i$ 为
\[
a_i^\dagger=z^i,\qquad a_i=\hbar\frac{\pa}{\pa z^i}.
\]
它们互为伴随，并满足
\[
[a_i,a_j^\dagger]=\delta_{ij}\hbar.
\]

% [原文 p.277]
给定多重指标 $I=\{i_1,i_2,\cdots,i_n\}$，我们定义多项式
\[
\lambda_I(\mathbf z):=\frac{\mathbf z^I}{\sqrt{\hbar^{|I|}I!}}
\]
其中
\[
\mathbf z^I:=(z^1)^{i_1}(z^2)^{i_2}\cdots(z^n)^{i_n},\qquad
|I|:=i_1+i_2+\cdots+i_n,\qquad
I!:=i_1!i_2!\cdots i_n!.
\]
若把 $1$ 认同为真空态，则
\[
\lambda_I=\frac{(a_1^\dagger)^{i_1}}{\sqrt{\hbar^{i_1}i_1!}}
\cdots\frac{(a_n^\dagger)^{i_n}}{\sqrt{\hbar^{i_n}i_n!}}1
\]
正是物理福克空间中激发态的表达式。

\begin{proposition}{5.3.7}
  $\{\lambda_I\}$ 构成 $\mathcal{F}_n$ 的一组正交归一基。
\end{proposition}

\begin{proof}
  正交性由 $a_i^\dagger$ 与 $a_i$ 的伴随关系得到。我们可以算出范数
  \[
  \begin{aligned}
  \langle\mathbf z^I|\mathbf z^I\rangle_\hbar
  &=\left\langle 1\left|a_1^{i_1}\cdots a_n^{i_n}
  (a_1^\dagger)^{i_1}\cdots(a_n^\dagger)^{i_n}\right|1\right\rangle_\hbar\\
  &=\hbar^{|I|}I!\langle1|1\rangle_\hbar=\hbar^{|I|}I!.
  \end{aligned}
  \]
  于是 $\{\lambda_I\}$ 是正交归一的集合。下面证明它构成一组基。
  对任意 $f\in\mathcal{F}_n$，考虑它的泰勒级数
  \[
  f(\mathbf z)=\sum_I c_I\mathbf z^I.
  \]
  该级数在 $\C^n$ 的紧子集上一致收敛到 $f$。设 $\chi_R$ 表示半径为 $R$ 的多圆盘上的示性函数
  \[
  \chi_R(\mathbf z)=\begin{cases}
  1 & \text{若 } \mathbf z\in P_R(0),\\
  0 & \text{其他}.
  \end{cases}
  \]
  用极坐标直接计算可得
  \[
  \langle\mathbf z^I\chi_R|\mathbf z^J\chi_R\rangle_\hbar=\delta_{IJ}d_{I,R},
  \]
  其中常数 $d_{I,R}$ 随 $R$ 增大而增大，且
  \[
  \lim_{R\to\infty}d_{I,R}=\hbar^{|I|}I!.
  \]
  在 $P_R(0)$ 上，级数的一致收敛性与上述正交性给出
  \[
  \norm{f\chi_R}_\hbar^2=\sum_I\abs{c_I}^2d_{I,R}.
  \]
  令 $R\to\infty$，得到
  \[
  \norm{f}_\hbar^2=\sum_I\abs{c_I}^2\norm{\mathbf z^I}_\hbar^2.
  \]
  这说明 $\{\lambda_I\}$ 构成一组基。
\end{proof}

% [原文 p.278]
\begin{proposition}{5.3.8}
  设 $f\in\mathcal{F}_n$。则
  \[
  \norm{z^if}_\hbar^2=\hbar\norm{f}_\hbar^2+\norm{\hbar\pa_{z^i}f}_\hbar^2.
  \]
  特别地，$a_i$ 与 $a_i^\dagger$ 在 $\mathcal{F}_n$ 上有相同的定义域。
\end{proposition}

\begin{proof}
  该恒等式对 $f=\lambda_I$ 成立，一般情形由帕塞瓦尔等式得到。
\end{proof}

下面考虑海森堡群 $H_n$ 在 Segal–Bargmann 空间 $\mathcal{F}_n$ 上的表示。回忆谐振子中升算子与降算子的构造
\[
\begin{cases}
a^\dagger=\frac{1}{\sqrt2}(\wh x-\ii\wh p)\\[4pt]
a=\frac{1}{\sqrt2}(\wh x+\ii\wh p)
\end{cases}
\]
或
\[
\begin{cases}
\wh x=\frac{1}{\sqrt2}(a+a^\dagger)\\[4pt]
\wh p=\frac{\ii}{\sqrt2}(a^\dagger-a)
\end{cases}
\]
设 $\R^{2n}$ 是标准辛空间，辛基为 $\{e_1,\cdots,e_n,f_1,\cdots,f_n\}$。回忆在薛定谔表示中，我们有
\[
e_i\longmapsto\frac{i}{\hbar}\wh x^i,\qquad
f_i\longmapsto-\frac{i}{\hbar}\wh p_i,\qquad
K\longmapsto\frac{i}{\hbar}.
\]
与上述升、降算子的公式比较，这提示我们在 $\mathcal{F}_n$ 上定义表示 $\rho$ 为
\[
e_i\longmapsto\frac{i}{\sqrt{2\hbar}}
\left(z^i+\hbar\frac{\pa}{\pa z^i}\right),
\]
\[
f_i\longmapsto\frac{1}{\sqrt{2\hbar}}
\left(z^i-\hbar\frac{\pa}{\pa z^i}\right),
\]
\[
K\longmapsto\frac{i}{\hbar}.
\]
于是
\[
\mathbf r\cdot\mathbf e+\mathbf s\cdot\mathbf f
\longmapsto\frac{1}{\sqrt{2\hbar}}(i\mathbf r+\mathbf s)\cdot\mathbf z
+\frac{1}{\sqrt2}(i\mathbf r-\mathbf s)\cdot\nabla_z.
\]
我们用复坐标 $\mathbf w\in\C^n$ 来表示 $(\mathbf r,\mathbf s)$：
\[
\begin{cases}
\mathbf w=\frac{1}{\sqrt2}(\mathbf s-i\mathbf r)\\[4pt]
\ol{\mathbf w}=\frac{1}{\sqrt2}(\mathbf s+i\mathbf r)
\end{cases}
\qquad\text{或}\qquad
\begin{cases}
\mathbf r=\frac{i}{\sqrt2}(\mathbf w-\ol{\mathbf w})\\[4pt]
\mathbf s=\frac{1}{\sqrt2}(\mathbf w+\ol{\mathbf w})
\end{cases}
\]
于是可以写成
\[
\mathbf r\cdot\mathbf e+\mathbf s\cdot\mathbf f
\longmapsto\frac{1}{\hbar}\ol{\mathbf w}\cdot\mathbf z-\mathbf w\cdot\nabla_z.
\]
这引导我们定义海森堡群的表示 $\rho$ 为
\[
\rho_\hbar(\mathbf w)=e^{\frac{1}{\hbar}\ol{\mathbf w}\cdot\mathbf z-\mathbf w\cdot\nabla_z}.
\]

% [原文 p.279]
用贝克–坎贝尔–豪斯多夫公式，这些算子满足复合法则
\[
\begin{aligned}
\rho_\hbar(\mathbf w_1)\rho_\hbar(\mathbf w_2)
&=\rho_\hbar(\mathbf w_1+\mathbf w_2)
e^{\frac{1}{2\hbar}(\ol{\mathbf w}_1\cdot\mathbf w_2-\mathbf w_1\cdot\ol{\mathbf w}_2)}\\
&=\rho_\hbar(\mathbf w_1+\mathbf w_2)
e^{\frac{i}{\hbar}\operatorname{Im}(\ol{\mathbf w}_1\cdot\mathbf w_2)}.
\end{aligned}
\]
特别地，我们有
\[
\rho_\hbar(\mathbf w)
=e^{\frac{1}{\hbar}\ol{\mathbf w}\cdot\mathbf z-\frac{1}{2\hbar}\abs{\mathbf w}^2}
e^{-\mathbf w\cdot\nabla_z}.
\]
于是我们得到 $\mathcal{F}_n$ 上表示的如下显式公式
\[
(\rho_\hbar(\mathbf w)f)(\mathbf z)
:=e^{\frac{1}{\hbar}\ol{\mathbf w}\cdot\mathbf z-\frac{1}{2\hbar}\abs{\mathbf w}^2}
f(\mathbf z-\mathbf w).
\]
这就是海森堡群的\textbf{巴格曼–福克表示}（Bargmann–Fock representation）。
由斯通–冯·诺依曼定理，存在酉映射
\[
\mathcal{B}:L^2(\R^n)\longrightarrow\mathcal{F}_n,
\]
它把薛定谔表示 $T_\hbar$ 与巴格曼–福克表示 $\rho_\hbar$ 交织起来
\[
\begin{array}{ccc}
L^2(\R^n) & \xrightarrow{\ \mathcal{B}\ } & \mathcal{F}_n\\[6pt]
\downarrow\ T_\hbar & & \downarrow\ \rho_\hbar\\[6pt]
L^2(\R^n) & \xrightarrow{\ \mathcal{B}\ } & \mathcal{F}_n
\end{array}
\]
显式地，这样的 $\mathcal{B}$ 由下式给出
\[
(\mathcal{B}\psi)(\mathbf z)
=\frac{1}{(\pi\hbar)^{n/4}}\int_{\R^n}\rd^n\mathbf x\,
e^{-\frac{1}{2\hbar}(\mathbf z\cdot\mathbf z-2\sqrt2\mathbf z\cdot\mathbf x+\mathbf x\cdot\mathbf x)}
\psi(\mathbf x).
\]
为理解这个公式，我们在李代数层面做核对。
\[
\begin{aligned}
\mathbf r\cdot\mathbf e+\mathbf s\cdot\mathbf f
&\xrightarrow{\ T_\hbar\ }\frac{i}{\hbar}\mathbf r\cdot\mathbf x-\mathbf s\cdot\nabla_x\\
&=-\frac{1}{\sqrt{2\hbar}}(\mathbf w-\ol{\mathbf w})\cdot\mathbf x
-\frac{1}{\sqrt2}(\mathbf w+\ol{\mathbf w})\cdot\nabla_x\\
&=\frac{1}{\hbar}\ol{\mathbf w}\cdot\left(\frac{1}{\sqrt2}\mathbf x
-\frac{\hbar}{\sqrt2}\nabla_x\right)
-\mathbf w\cdot\left(\frac{1}{\sqrt{2\hbar}}\mathbf x
+\frac{1}{\sqrt2}\nabla_x\right),
\end{aligned}
\]
\[
\mathbf r\cdot\mathbf e+\mathbf s\cdot\mathbf f
\xrightarrow{\ \rho_\hbar\ }
\frac{1}{\hbar}\ol{\mathbf w}\cdot\mathbf z-\mathbf w\cdot\nabla_z.
\]
比较上面两个无穷小变换，$\mathcal{B}$ 应当具有性质
\[
z^i(\mathcal{B}\psi)
=\mathcal{B}\left[\left(\frac{1}{\sqrt2}x^i
-\frac{\hbar}{\sqrt2}\pa_{x^i}\right)\psi\right],
\]
\[
\pa_{z^i}(\mathcal{B}\psi)
=\mathcal{B}\left[\left(\frac{1}{\sqrt{2\hbar}}x^i
+\frac{1}{\sqrt2}\pa_{x^i}\right)\psi\right].
\]
不难验证，上面定义的 $\mathcal{B}$ 恰好满足这两个等式。

% [原文 p.280]
我们再核对酉性。设 $\psi\in\mathcal{S}(\R^n)$ 是施瓦茨函数。则
\[
\begin{aligned}
\norm{\mathcal{B}\psi}_\hbar^2
&=\frac{1}{(\pi\hbar)^{3n/2}}\int\rd^n\mathbf x\rd^n\mathbf y\,
\psi(\mathbf x)\ol{\psi(\mathbf y)}\\
&\qquad\int\rd^{2n}\mathbf z\,
e^{-\frac{1}{2\hbar}\left(\mathbf z\cdot\mathbf z-2\sqrt2\mathbf z\cdot\mathbf x
+\mathbf x\cdot\mathbf x+\ol{\mathbf z}\cdot\ol{\mathbf z}
-2\sqrt2\ol{\mathbf z}\cdot\mathbf y+\mathbf y\cdot\mathbf y\right)}
e^{-\abs{\mathbf z}^2/\hbar}\\
&\stackrel{\mathbf z=(\mathbf r+i\mathbf s)/\sqrt2}{=}
\frac{1}{2^n(\pi\hbar)^{3n/2}}\int\rd^n\mathbf x\rd^n\mathbf y\,
\psi(\mathbf x)\ol{\psi(\mathbf y)}\\
&\qquad\int\rd^n\mathbf r\rd^n\mathbf s\,
e^{-\frac{1}{\hbar}\left(\mathbf r^2-\mathbf r\cdot(\mathbf x+\mathbf y)
+i\mathbf s\cdot(\mathbf x-\mathbf y)+\frac12(\mathbf x\cdot\mathbf x+\mathbf y\cdot\mathbf y)\right)}\\
&=\frac{1}{(2\pi\hbar)^n}\int\rd^n\mathbf x\rd^n\mathbf y\,
\psi(\mathbf x)\ol{\psi(\mathbf y)}
e^{-\frac{1}{4\hbar}(\mathbf x\cdot\mathbf x-\mathbf y\cdot\mathbf y)^2}
\int\rd^n\mathbf s\,e^{-\frac{i}{\hbar}\mathbf s\cdot(\mathbf x-\mathbf y)}\\
&=\int\rd^n\mathbf x\rd^n\mathbf y\,\psi(\mathbf x)\ol{\psi(\mathbf y)}
e^{-\frac{1}{4\hbar}(\mathbf x\cdot\mathbf x-\mathbf y\cdot\mathbf y)^2}
\delta(\mathbf x-\mathbf y)\\
&=\int\rd^n\mathbf x\,\psi(\mathbf x)\ol{\psi(\mathbf x)}
=\norm{\psi}_{L^2(\R^n)}^2,
\end{aligned}
\]
\tutornote{第三、四行的指数按原文录入，依上一步推导应为 $-\frac{1}{4\hbar}(\mathbf x-\mathbf y)\cdot(\mathbf x-\mathbf y)$。}
这与预期相符。这个酉映射 $\mathcal{B}:L^2(\R^n)\to\mathcal{F}_n$ 称为
\textbf{Segal–Bargmann 变换}。

\takeaway{本节从辛向量空间出发构造海森堡李代数与海森堡群，给出薛定谔表示与巴格曼–福克表示，并用斯通–冯·诺依曼定理说明二者由 Segal–Bargmann 变换酉等价。}

\section{几何量子化}
\subsection{狄拉克量子化原理}

经典力学植根于相空间流形 $M$ 的辛几何。经典可观测量的空间
$\mathrm{Obs}^{cl}:=C^\infty(M)$ 带有泊松括号 $\{-,-\}$ 的结构。
若对 $g\in\mathrm{Obs}^{cl}$，
\[
\{g,f_1\}=\cdots=\{g,f_m\}=0
\quad\Longrightarrow\quad g=\text{常数},
\]
则称 $\{f_1,\cdots,f_m\}$ 为经典可观测量的一个完全集。
例如在 $M=\R^{2n}$ 上，位置与动量函数构成的集合
$\{x^1,\cdots,x^n,p_1,\cdots,p_n\}$ 是完备的。

量子力学中，我们有表示量子态空间的希尔伯特空间 $\mathbb{H}$。量子可观测量的空间
$\mathrm{Obs}^{q}$ 是 $\mathbb{H}$ 上适当的线性算子族，它带有算子交换子的自然结构
\[
[\mathcal{O}_1,\mathcal{O}_2]
=\mathcal{O}_1\mathcal{O}_2-\mathcal{O}_2\mathcal{O}_1.
\]
类似地，若对 $\mathcal{O}\in\mathrm{Obs}^{q}$，
\[
[\mathcal{O},\mathcal{O}_1]=\cdots=[\mathcal{O},\mathcal{O}_m]=0
\quad\Longrightarrow\quad
\mathcal{O}=c\mathbf{1}\ \text{是单位算子的倍数},
\]
则称量子可观测量 $\{\mathcal{O}_1,\cdots,\mathcal{O}_m\}$ 是量子可观测量的一个完全集。

在狄拉克原理（Dirac's Principle）中，量子化归结为给每个经典观测量 $f$ 指定某个希尔伯特空间上的算子 $\wh f$：
\[
\begin{aligned}
Q&:\mathrm{Obs}^{cl}\longrightarrow\mathrm{Obs}^{q}\\
f&\longmapsto\wh f.
\end{aligned}
\]
这个指定必须满足

Q1：$Q$ 是 $\R$ 线性的。

Q2：$\wh 1=Q(1)$ 是单位算子。

% [原文 p.281]
Q3：若 $f$ 是实值函数，则 $\wh f$ 自伴。

Q4：
\[
[\wh f,\wh g]=i\hbar\widehat{\{f,g\}}.
\]

Q5：若 $\{f_1,\cdots,f_m\}$ 是经典可观测量的完全集，则
$\{\wh f_1,\cdots,\wh f_m\}$ 是量子可观测量的完全集。

Q1、Q2、Q3 是自然的要求。Q4 是狄拉克的观察：算子交换子应当是经典泊松括号的量子对应物。Q5 可以看作一个不可约性条件。希尔伯特空间 $\mathbb{H}$ 一般是无限维的，但也可能取有限维（例如描述内部对称性时）。
不幸的是，一般不可能同时满足 Q4 与 Q5。这一点可由量子化的一个“不可行”（no-go）型结果说明，我们将在 5.4.2 节讨论。
在此之前，我们先看薛定谔表示上的外尔量子化。记外尔量子化（见定义 3.2.7）为
\[
\mathcal{W}_\hbar:f\mapsto f^W.
\]
设 $f(\mathbf x,\mathbf p)$ 是多项式。它的外尔量子化是算子
\[
\mathcal{W}_\hbar(f)=\mathrm{Sym}(f(\wh x^i,\wh p_i)),
\]
其中
\[
\wh x^i=x^i,\qquad \wh p_i=-i\hbar\frac{\pa}{\pa x^i},
\]
而 $\mathrm{Sym}(-)$ 是对称序平均。我们有
\[
\mathcal{W}_\hbar(f)\mathcal{W}_\hbar(g)=\mathcal{W}_\hbar(f*_\hbar g),
\]
其中 $f*_\hbar g$ 是 Moyal 乘积
\[
(f*_\hbar g)(\mathbf x,\mathbf p)
=f(\mathbf x,\mathbf p)
e^{\frac{i\hbar}{2}\left(\frac{\overleftarrow{\pa}}{\pa x^i}
\frac{\overrightarrow{\pa}}{\pa p_i}
-\frac{\overleftarrow{\pa}}{\pa p_i}
\frac{\overrightarrow{\pa}}{\pa x^i}\right)}
g(\mathbf x,\mathbf p).
\]
特别地，
\[
[\mathcal{W}_\hbar(f),\mathcal{W}_\hbar(g)]
=\mathcal{W}_\hbar([f,g]_*),
\]
其中
\[
[f,g]_*:=f*_\hbar g-g*_\hbar f
\]
是 Moyal 乘积的交换子。计算得
\[
[f,g]_*=i\hbar\{f,g\}
+\frac{(i\hbar)^3}{24}f\left(\frac{\overleftarrow{\pa}}{\pa x^i}
\frac{\overrightarrow{\pa}}{\pa p_i}
-\frac{\overleftarrow{\pa}}{\pa p_i}
\frac{\overrightarrow{\pa}}{\pa x^i}\right)^3g+\cdots.
\]
因此在薛定谔表示中，外尔量子化违反 Q4。
另一方面，上述计算表明：若 $f,g$ 是次数不超过 $2$ 的多项式，则
\[
[f,g]_*=i\hbar\{f,g\}.
\]
我们随后会看到，次数为 $4$ 时就会出现障碍。

% [原文 p.282]
\subsection{格罗嫩沃尔德–范霍夫定理}

格罗嫩沃尔德–范霍夫定理（Groenewold–Van Hove Theorem）给出了量子化的若干“不可行”（No-Go）型结果。下面叙述其中最简单的版本。
考虑 $\R^{2n}$ 上的量子化。设
\[
P^{\le k}:=\{f\in\R[x^i,p_i]\mid \deg f\le k\}.
\]
薛定谔表示把 $P^{\le1}$ 中的元素对应到算子
\[
1\longmapsto\wh 1=1,
\]
\[
x^i\longmapsto\wh x^i=x^i,
\]
\[
p_i\longmapsto\wh p_i=-i\hbar\frac{\pa}{\pa x^i},
\]
并且
\[
[\wh x^i,\wh p_j]=i\hbar\delta_{ij}
=i\hbar\widehat{\{x^i,p_j\}}.
\]

\begin{theorem}{5.4.1}
  （Groenewold）不存在 $\R$ 线性映射
  \[
  Q:P^{\le4}\longrightarrow L^2(\R^n)\text{ 上的自伴算子},
  \]
  使得
  \[
  (1)\quad Q(1)=1,\quad Q(x^i)=\wh x^i,\quad Q(p_i)=\wh p_i,
  \]
  \[
  (2)\quad [Q(f),Q(g)]=i\hbar Q(\{f,g\}),\qquad
  \forall f,g\in P^{\le4}.
  \]
\end{theorem}

\begin{proof}
  我们证明 $n=1$ 的情形，$n>1$ 的情形的证明类似。假设这样的 $Q$ 存在。沿用上面的记号，用 $\mathcal{W}_\hbar:f\mapsto f^W$ 表示外尔量子化。由假设 (1)，
  \[
  Q(f)=\mathcal{W}_\hbar(f)\qquad\text{对 } f\in P^{\le1}.
  \]
  注意泊松括号满足
  \[
  \{P^{\le2},P^{\le1}\}\subset P^{\le1},
  \]
  并且
  \[
  [\mathcal{W}_\hbar(f),\mathcal{W}_\hbar(g)]
  =i\hbar\mathcal{W}_\hbar(\{f,g\})
  \qquad\text{对 } f\in P^{\le2},\ g\in P^{\le1}.
  \]
  因此
  \[
  [\mathcal{W}_\hbar(f),\wh g]
  =i\hbar\widehat{\{f,g\}}
  =[Q(f),\wh g]
  \qquad\text{对 } f\in P^{\le2},\ g\in P^{\le1}.
  \]
  于是
  \[
  [Q(f)-\mathcal{W}_\hbar(f),\wh x]
  =[Q(f)-\mathcal{W}_\hbar(f),\wh p]=0.
  \]
  由此可知，$L^2(\R^n)$ 上的单参数酉算子族
  $e^{i\alpha(Q(f)-\mathcal{W}_\hbar(f))}$ 与海森堡群作用交换，并给出薛定谔表示的自同构。

% [原文 p.283]
  由斯通–冯·诺依曼定理（定理 5.3.3），这类自同构必为标量乘法，因此
  \[
  Q(f)-\mathcal{W}_\hbar(f)=\text{const}\ 1
  \qquad\text{对 } f\in P^{\le2}.
  \]
  右端的这些常数实际上是零。事实上，利用 $x^2=\frac12\{x^2,xp\}$，
  \[
  \begin{aligned}
  Q(x^2)&=\frac{1}{2i\hbar}[Q(x^2),Q(xp)]\\
  &=\frac{1}{2i\hbar}[\mathcal{W}_\hbar(x^2)+c_1,\mathcal{W}_\hbar(xp)+c_2]
  \qquad c_1,c_2\text{ 为某些常数}\\
  &=\frac{1}{2i\hbar}[\mathcal{W}_\hbar(x^2),\mathcal{W}_\hbar(xp)]\\
  &=\mathcal{W}_\hbar(x^2).
  \end{aligned}
  \]
  类似地可以证明 $Q(p^2)=\mathcal{W}_\hbar(p^2)$ 与 $Q(xp)=\mathcal{W}_\hbar(xp)$。因此
  \[
  Q(f)=\mathcal{W}_\hbar(f)\qquad\text{对 } f\in P^{\le2}.
  \]
  下面注意，对 $f\in P^{\le3}$、$g\in P^{\le1}$，
  \[
  [\mathcal{W}_\hbar(f),\mathcal{W}_\hbar(g)]
  =i\hbar\mathcal{W}_\hbar(\{f,g\})
  \]
  仍然成立。于是对 $f\in P^{\le3}$、$g\in P^{\le1}$，
  \[
  [Q(f)-\mathcal{W}_\hbar(f),\wh g]
  =i\hbar Q(\{f,g\})-i\hbar\mathcal{W}_\hbar(\{f,g\})=0,
  \]
  因为 $\{f,g\}\in P^{\le2}$。用与上面相同的论证，我们有
  \[
  Q(f)-\mathcal{W}_\hbar(f)=\text{const}\ 1
  \qquad\text{对 } f\in P^{\le3}.
  \]
  类似地，利用
  \[
  x^3=\frac13\{x^3,xp\},\qquad p^3=\frac13\{xp,p^3\},
  \]
  \[
  x^2p=\frac16\{x^3,p^2\},\qquad xp^2=\frac16\{x^2,p^3\},
  \]
  可以推出右端的这些常数都是零。因此
  \[
  Q(f)=\mathcal{W}_\hbar(f)\qquad\text{对 } f\in P^{\le3}.
  \]
  现在来看次数为 $4$ 的多项式。考虑
  \[
  x^2p^2=\frac19\{x^3,p^3\}=\frac13\{x^2p,xp^2\}.
  \]
  由假设 (2)，我们应当有
  \[
  Q(x^2p^2)=\frac{1}{9i\hbar}[Q(x^3),Q(p^3)]
  =\frac{1}{3i\hbar}[Q(x^2p),Q(xp^2)].
  \]
  另一方面，利用
  \[
  [x^3,p^3]_*=9(i\hbar)x^2p^2+\frac32(i\hbar)^3,
  \]
  \[
  [x^2p,xp^2]_*=3(i\hbar)x^2p^2-\frac12(i\hbar)^3,
  \]

% [原文 p.284]
  我们有
  \[
  \begin{aligned}
  \frac{1}{9i\hbar}[Q(x^3),Q(p^3)]
  &=\frac{1}{9i\hbar}[\mathcal{W}_\hbar(x^3),\mathcal{W}_\hbar(p^3)]\\
  &=\frac{1}{9i\hbar}\mathcal{W}_\hbar([x^3,p^3]_*)
  =\mathcal{W}_\hbar(x^2p^2)-\frac16\hbar^2
  \end{aligned}
  \]
  以及
  \[
  \begin{aligned}
  \frac{1}{3i\hbar}[Q(x^2p),Q(xp^2)]
  &=\frac{1}{3i\hbar}[\mathcal{W}_\hbar(x^2p),\mathcal{W}_\hbar(xp^2)]\\
  &=\frac{1}{3i\hbar}\mathcal{W}_\hbar([x^2p,xp^2]_*)
  =\mathcal{W}_\hbar(x^2p^2)+\frac16\hbar^2\\
  &\ne\frac{1}{9i\hbar}[Q(x^3),Q(p^3)].
  \end{aligned}
  \]
  矛盾。
\end{proof}

\subsection{预量子化}

由于构造满足全部条件 Q1–Q5 的量子化并不现实，我们必须放宽条件 Q4 或 Q5。有一种自然的方式为辛流形构造满足条件 Q1–Q4 的量子化，这就是\textbf{预量子化}（prequantization）。
设 $(M,\omega)$ 是辛流形。回忆哈密顿向量场的构造
\[
\begin{aligned}
C^\infty(M)&\longrightarrow\mathrm{Vect}(M)\\
f&\longmapsto V_f,\qquad \iota_{V_f}\omega=\rd f,
\end{aligned}
\]
它是李代数同态
\[
[V_f,V_g]=V_{\{f,g\}}.
\]
若把 $V_f$ 看作作用在 $C^\infty(M)$ 上的一阶微分算子，则对应
\[
f\longmapsto i\hbar V_f
\]
满足 Q1、Q3、Q4，但不满足 Q2，因为它把 $1$ 送到 $0$。
为弥补这一点，设 $L\to M$ 是 $M$ 上带联络 $\nabla$ 的复线丛。设 $F_\nabla$ 是曲率 2-形式。对任意向量场 $X,Y$，我们有
\[
[\nabla_X,\nabla_Y]-\nabla_{[X,Y]}=F_\nabla(X,Y).
\]
给定函数 $f$，我们用 $\nabla_f$ 表示沿哈密顿向量场 $V_f$ 的协变导数
\[
\nabla_f:=\nabla_{V_f}.
\]
于是上面的曲率公式成为
\[
[\nabla_f,\nabla_g]-\nabla_{\{f,g\}}
=F_\nabla(V_f,V_g)=\iota_{V_g}\iota_{V_f}F_\nabla.
\]
这里 $\iota_{V_f}$ 是沿向量场 $V_f$ 的内乘（缩并）。
考虑在 $\Gamma(M,L)$ 上如下指定算子
\[
Q(f):=i\hbar\nabla_f+f.
\]

% [原文 p.285]
我们有
\[
\begin{aligned}
[Q(f),Q(g)]
&=[i\hbar\nabla_f+f,i\hbar\nabla_g+g]\\
&=(i\hbar)^2\left(\nabla_{\{f,g\}}+F_\nabla(V_f,V_g)\right)
+2i\hbar\{f,g\}\\
&=(i\hbar)\left(i\hbar\nabla_{\{f,g\}}+\{f,g\}\right)
+(i\hbar)^2F_\nabla(V_f,V_g)+i\hbar\{f,g\}\\
&=i\hbar Q(\{f,g\})+(i\hbar)^2F_\nabla(V_f,V_g)+i\hbar\{f,g\}\\
&=i\hbar Q(\{f,g\})+(i\hbar)^2F_\nabla(V_f,V_g)
+i\hbar\iota_{V_f}\iota_{V_g}\omega.
\end{aligned}
\]
于是，若曲率 2-形式为
\[
F_\nabla=-\frac{i}{\hbar}\omega,
\]
则下列关系成立
\[
[Q(f),Q(g)]=i\hbar Q(\{f,g\}).
\]
这个计算引出下面的定义。

\begin{definition}{5.4.2}
  辛流形 $(M,\omega)$ 上的\textbf{预量子化线丛}（prequantum line bundle）是指厄米线丛
  $(\mathcal{L},\langle-,-\rangle)$，连同相容的联络 $\nabla$，使得其曲率 2-形式
  $F_\nabla$ 满足
  \[
  F_\nabla=-\frac{i}{\hbar}\omega.
  \]
\end{definition}

这里回忆一下，厄米线丛是复线丛 $\mathcal{L}$，在其纤维上带有光滑变化的厄米内积族
$\langle-,-\rangle$。我们对厄米内积的约定是：$\langle-,-\rangle$ 关于第一个变量共轭线性，关于第二个变量复线性。于是对任意局部截面 $s,s_1,s_2\in\Gamma(U,\mathcal{L})$ 与局部函数 $f\in C^\infty(U)$，
\begin{itemize}
  \item $\ol{\langle s_1,s_2\rangle}=\langle s_2,s_1\rangle$；
  \item $\langle fs_1,s_2\rangle=\ol f\langle s_1,s_2\rangle$；
  \item $\langle s_1,fs_2\rangle=f\langle s_1,s_2\rangle$；
  \item $\langle s,s\rangle\ge0$，且 $\langle s,s\rangle=0\iff s=0$。
\end{itemize}
称联络 $\nabla$ 与厄米内积相容，是指对 $M$ 上任意实向量场 $X$ 与 $\mathcal{L}$ 的局部截面 $s_1,s_2$，
\[
X\langle s_1,s_2\rangle
=\langle\nabla_Xs_1,s_2\rangle+\langle s_1,\nabla_Xs_2\rangle.
\]
厄米结构 $\langle-,-\rangle$ 在 $\mathcal{L}$ 的整体截面之间诱导厄米内积 $\langle-|-\rangle$：
\[
\langle s_1|s_2\rangle
:=\int_M\langle s_1,s_2\rangle\frac{\omega^n}{n!}.
\]
取 $\mathcal{L}$ 的平方可积可测截面，我们得到希尔伯特空间 $L^2(M,\mathcal{L})$。于是
\[
\norm{s}:=\left(\int_M\langle s,s\rangle\frac{\omega^n}{n!}\right)^{\frac12}<\infty
\qquad\text{对 } s\in L^2(M,\mathcal{L}).
\]

% [原文 p.286]
设 $(\mathcal{L},\omega,\langle-,-\rangle)$ 是 $(M,\omega)$ 上的预量子化线丛。对光滑函数 $f$，我们指定如下算子
\[
Q(f)=i\hbar\nabla_f+f:\
\Gamma(M,\mathcal{L})\longrightarrow\Gamma(M,\mathcal{L}).
\]
曲率条件 $F_\nabla=-\frac{i}{\hbar}\omega$ 推出
\[
[Q(f),Q(g)]=i\hbar Q(\{f,g\}),
\]
因此条件 Q4 成立。此外我们有
\[
Q(1)=1,
\]
因此条件 Q2 成立。关于 $\R$ 线性的条件 Q1 是显然的。
下面考察条件 Q3。设 $s_1,s_2\in\Gamma_c(M,\mathcal{L})$ 是 $\mathcal{L}$ 的两个具紧支集的截面，$f$ 是 $M$ 上的实光滑函数。$\nabla$ 与 $\langle-,-\rangle$ 的相容性给出
\[
\nabla_f\langle s_1,s_2\rangle
=\langle\nabla_fs_1,s_2\rangle+\langle s_1,\nabla_fs_2\rangle.
\]
按体积形式 $\frac{\omega^n}{n!}$ 在 $M$ 上积分，得到
\[
\begin{aligned}
\langle\nabla_fs_1|s_2\rangle+\langle s_1|\nabla_fs_2\rangle
&=\int_M\left(\langle\nabla_fs_1,s_2\rangle+\langle s_1,\nabla_fs_2\rangle\right)
\frac{\omega^n}{n!}\\
&=\int_M\nabla_f\langle s_1,s_2\rangle\,\frac{\omega^n}{n!}\\
&=\int_M\left(\mathscr{L}_{V_f}\langle s_1,s_2\rangle\right)
\frac{\omega^n}{n!}\\
&=-\int_M\langle s_1,s_2\rangle\,
\mathscr{L}_{V_f}\left(\frac{\omega^n}{n!}\right)=0.
\end{aligned}
\]
这里最后一步用了刘维尔定理：刘维尔体积形式在哈密顿流下不变。这说明自伴关系
\[
\langle Q(f)s_1|s_2\rangle=\langle s_1|Q(f)s_2\rangle.
\]
于是预量子化线丛给出满足条件 Q1–Q4 的量子化。
上面的讨论以预先给定预量子化线丛为前提。然而预量子化线丛是否存在并不清楚。它确实存在拓扑上的障碍。
2-形式
\[
\frac{i}{2\pi}F_\nabla
\]
是闭的，其德拉姆上同调类
\[
\left[\frac{i}{2\pi}F_\nabla\right]=c_1(\mathcal{L})
\]
表示 $\mathcal{L}$ 的第一陈类。陈类的一个显著性质是整性：
\[
\frac{i}{2\pi}\int_\Sigma F_\nabla\in\Z
\]
对任意闭曲面 $\Sigma$ 成立。因此 $\frac{1}{2\pi\hbar}[\omega]$ 表示一个整上同调类。反过来也对：若 $\frac{1}{2\pi\hbar}[\omega]$ 是整的，则存在预量子化线丛。

% [原文 p.287]
\begin{definition}{5.4.3}
  若对 $M$ 中每个闭曲面 $\Sigma$ 都有
  \[
  \frac{1}{2\pi\hbar}\int_\Sigma\omega\in\Z,
  \]
  则称辛流形 $(M,\omega)$ 可量子化（quantizable，对 $\hbar$ 的某个特定取值而言）。
\end{definition}

于是在可量子化的相空间上，我们能构造满足 Q1–Q4 的量子化。注意这种量子化与量子力学中的预期相去甚远。例如
\[
Q(f)=i\hbar\nabla_f+f
\]
对任意函数 $f$ 始终是一阶微分算子。但在量子力学中，我们期待出现更高阶的微分算子，例如薛定谔算子。问题出在条件 Q5：上面构造的希尔伯特空间 $L^2(M,\mathcal{L})$ 太大。我们将在 5.4.4 节讨论如何削减自由度。

\begin{example}{5.4.4}
  考虑余切丛的全空间
  \[
  M=T^*X,\qquad \omega=d\theta,
  \]
  其中 $\theta$ 是刘维尔 1-形式。由于 $\omega$ 是恰当的，
  \[
  \int_\Sigma\omega=\int_\Sigma d\theta=0
  \]
  对任意闭曲面 $\Sigma\subset M$ 成立。因此余切丛总是可量子化的。事实上，我们可以取 $\mathcal{L}$ 为平凡线丛，并取联络为
  \[
  \nabla=d-\frac{i}{\hbar}\theta.
  \]
\end{example}

\begin{example}{5.4.5}
  $M=S^2=\{x^2+y^2+z^2=1\}\subset\R^3$，
  \[
  \omega=A(x\rd y\wedge\rd z+y\rd z\wedge\rd x+z\rd x\wedge\rd y),
  \qquad A>0.
  \]
  我们有
  \[
  \int_{S^2}\omega=\int_{x^2+y^2+z^2\le1}d\omega=4\pi A.
  \]
  于是 $(M,\omega)$ 可量子化 $\iff A\in\frac{\hbar}{2}\Z$。
\end{example}

\subsection{极化与量子态}

由可量子化辛流形 $M$ 上预量子化线丛 $\mathcal{L}$ 得到的希尔伯特空间
$L^2(M,\mathcal{L})$，问题在于它太大。例如在 $M=\R^{2n}$ 上，$\mathcal{L}$ 的截面 $\psi$ 同时依赖于位置 $\mathbf x$ 与动量 $\mathbf p$，因此不能充当量子力学中波函数的候选者。我们需要一种办法，消去“一半”的变量。极化的概念正是为此而设。

% [原文 p.288]
设 $T^\C M=TM\otimes_{\R}\C$ 表示切丛 $TM$ 的复化。设
\[
\bar{\ }:T^\C M\longrightarrow T^\C M,\qquad X\longmapsto\ol X
\]
表示复共轭。辛形式 $\omega$ 复线性地延拓为 $T^\C M$ 上的非退化配对，我们仍记之为 $\omega$。
设 $\dim_{\R}(M)=2n$。对任意 $z\in M$，若 $\C$ 线性子空间 $P\subset T_z^\C M$ 满足 $\dim_{\C}P=n$ 且
\[
\omega(X,Y)=0,\qquad\forall X,Y\in P,
\]
则称 $P$ 为拉格朗日子空间。

\begin{definition}{5.4.6}
  $(M,\omega)$ 的\textbf{极化}（polarization）是指复子丛 $P\subset T^\C M$，满足
  
  ① 对任意 $z\in M$，$P_z$ 是 $T_z^\C M$ 的拉格朗日子空间；
  
  ② 可积性：对任意开子集 $U$，
  \[
  X,Y\in\Gamma(U,P)\quad\Longrightarrow\quad[X,Y]\in\Gamma(U,P);
  \]
  这里 $[X,Y]$ 是两个向量场的李括号。
  
  ③ $\dim(P_z\cap\ol P_z)$ 关于 $z\in M$ 为常数。
\end{definition}

条件 ③ 是为某些构造而作的技术性假设。
我们主要关心下面两类极化：
\begin{itemize}
  \item[a)] \textbf{实极化}（Real Polarization）：$P_z=\ol P_z$，$\forall z\in M$；
  \item[b)] \textbf{凯勒极化}（Kähler Polarization）：$P_z\cap\ol P_z=\{0\}$，$\forall z\in M$。
\end{itemize}

\begin{definition}{5.4.7}
  设 $(\mathcal{L},\nabla)$ 是 $(M,\omega)$ 上的预量子化线丛，$P$ 是一个极化。若截面
  $s\in\Gamma(M,\mathcal{L})$ 满足
  \[
  \nabla_Xs=0
  \]
  对 $P$ 中任意向量场 $X$ 成立，则称 $s$ 是 $P$-极化的。
\end{definition}

我们来检验上述条件的相容性。设 $X,Y$ 是 $P$ 中的两个向量场。假设 $\nabla_Xs=\nabla_Ys=0$ 成立，则
\[
[\nabla_X,\nabla_Y]s=0.
\]
另一方面，
\[
\begin{aligned}
[\nabla_X,\nabla_Y]
&=\nabla_{[X,Y]}+F_\nabla(X,Y)\\
&=\nabla_{[X,Y]}-\frac{i}{\hbar}\omega(X,Y).
\end{aligned}
\]
由于 $P$ 是极化，$[X,Y]$ 仍是 $P$ 中的向量，且 $\omega(X,Y)=0$。这说明极化截面的定义方程是相容的。

% [原文 p.289]
于是我们可以定义与预量子化线丛 $\mathcal{L}$ 及极化 $P$ 相伴的量子希尔伯特空间
\[
L^2(M,\mathcal{L})_P
:=\{s\in L^2(M,\mathcal{L})\mid s\text{ 是 }P\text{-极化的}\}.
\]

\begin{example}{5.4.8}
  考虑相空间 $\R^{2n}$，配以如下预量子化线丛：
  \begin{itemize}
    \item $\mathcal{L}$ 是 $\R^{2n}$ 上的平凡线丛，其厄米内积为
    \[
    \langle f,g\rangle=\ol f g;
    \]
    \item 联络 $\nabla$ 为
    \[
    \nabla=d-\frac{i}{\hbar}\theta,
    \]
    其中 $\theta=\sum_i p_i\rd x^i$ 是刘维尔 1-形式。$\nabla$ 的曲率
    \[
    F_\nabla=-\frac{i}{\hbar}d\theta=-\frac{i}{\hbar}\omega
    \]
    满足预量子化条件。
  \end{itemize}
  我们考虑由沿动量方向的向量场张成的极化
  \[
  P=\Span\left\{\frac{\pa}{\pa p_i}\right\}.
  \]
  对任意函数 $f(\mathbf x,\mathbf p)$，
  \[
  \nabla_{\frac{\pa}{\pa p_i}}f=\frac{\pa f}{\pa p_i}
  -\frac{i}{\hbar}\theta\left(\frac{\pa}{\pa p_i}\right)f
  =\frac{\pa f}{\pa p_i}.
  \]
  因此 $f$ 是 $P$-极化的，当且仅当
  \[
  f=f(\mathbf x)
  \]
  不依赖于 $\mathbf p$。这给出了波函数的预期结果。
  若考虑由沿位置方向的向量场张成的另一极化
  \[
  P'=\Span\left\{\frac{\pa}{\pa x^i}\right\},
  \]
  则由
  \[
  \nabla_{\frac{\pa}{\pa x^i}}f=\frac{\pa f}{\pa x^i}
  -\frac{i}{\hbar}\theta\left(\frac{\pa}{\pa x^i}\right)f
  =\frac{\pa f}{\pa x^i}-\frac{i}{\hbar}p_if,
  \]
  知 $f$ 是 $P'$-极化的，当且仅当 $f$ 形如
  \[
  f=e^{\frac{i}{\hbar}\mathbf x\cdot\mathbf p}g(\mathbf p).
  \]
\end{example}

\begin{example}{5.4.9}
  考虑相空间 $\C^n$，配以凯勒辛形式
  \[
  \omega=i\pa\ol\pa K.
  \]

  % [原文 p.290]
  这里 $K$ 是实函数，称为凯勒势。我们取预量子化线丛 $\mathcal{L}$ 为平凡线丛，其厄米内积为
  \[
  \langle f_1,f_2\rangle=\ol f_1f_2e^{-\frac{K}{\hbar}},
  \]
  相容联络为
  \[
  \nabla=d-\frac{1}{\hbar}\pa K.
  \]
  其曲率
  \[
  F_\nabla=-\frac{1}{\hbar}d(\pa K)
  =-\frac{1}{\hbar}\ol\pa\pa K
  =\frac{1}{\hbar}\pa\ol\pa K
  =-\frac{i}{\hbar}\omega
  \]
  满足预量子化条件。
  我们取由 $(0,1)$ 型向量场张成的凯勒极化
  \[
  P=\Span\left\{\frac{\pa}{\pa\ol z^1},\cdots,\frac{\pa}{\pa\ol z^n}\right\}.
  \]
  由于
  \[
  \nabla_{\frac{\pa}{\pa\ol z^i}}f=\frac{\pa f}{\pa\ol z^i},
  \]
  $f$ 是 $P$-极化的，当且仅当 $f=f(\mathbf z)$ 全纯。由此得到的量子希尔伯特空间由满足
  \[
  \int_{\C^n}\abs{f(\mathbf z)}^2e^{-\frac{K}{\hbar}}\frac{\omega^n}{n!}<\infty
  \]
  的全纯函数 $f(\mathbf z)$ 组成。
  取 $K=\abs{\mathbf z}^2$，我们恰好得到 Segal–Bargmann 空间。
\end{example}

\begin{example}{5.4.10}
  我们把上述构造推广到带厄米全纯线丛 $\mathcal{L}$ 的复流形 $M$。设 $\nabla$ 是 $\mathcal{L}$ 上的相容联络，使得
  \[
  \nabla=\nabla^{1,0}+\nabla^{0,1},\qquad\text{其中 } \nabla^{0,1}=\ol\pa.
  \]
  假设第一陈类 $\frac{i}{2\pi}F_\nabla$ 在 $M$ 上定义辛形式。
  我们考虑相空间 $(M,\omega)$，其中
  \[
  \omega=i\hbar F_\nabla.
  \]
  按构造，$(\mathcal{L},\nabla)$ 给出预量子化线丛。我们取极化
  \[
  P=T^{0,1}M,
  \]
  它由 $(0,1)$ 型向量场张成。在局部全纯坐标 $\{z^i\}$ 下，$P$ 由
  $\left\{\frac{\pa}{\pa\ol z^i}\right\}$ 张成。因此 $P$-极化截面就是 $\mathcal{L}$ 的全纯截面。
  设 $M$ 紧。则全纯截面的空间 $H^0(M,\mathcal{L})$ 有限维。于是我们得到有限维的量子希尔伯特空间
  \[
  \mathbb{H}=H^0(M,\mathcal{L}).
  \]
  这样的希尔伯特空间可以用来描述量子系统的内部自由度。

% [原文 p.291]
  作为例子，考虑
  \[
  M=\CP^1=\{\C^2\text{ 中的直线}\}.
  \]
  它由两张全纯图卡覆盖
  \[
  z\in U=\C,\qquad w\in V=\C,
  \]
  它们在 $\C^*$ 上通过认同
  \[
  z=\frac1w
  \]
  粘合。我们考虑全纯线丛
  \[
  \mathcal{L}=\mathcal{O}_{\CP^1}(k),\qquad k\in\Z^{>0},
  \]
  其从 $U$ 到 $V$ 的转移函数在交叠 $U\cap V$ 上为 $w^k$。全纯截面由一对 $(f(z),g(w))$ 给出，其中
  \begin{itemize}
    \item $f(z)$ 在 $U$ 上全纯；
    \item $g(w)$ 在 $V$ 上全纯；
    \item 在 $U\cap V$ 上 $f(z)w^k=g(w)$。
  \end{itemize}
  我们发现 $f(z)$ 只能是次数 $\le k$ 的多项式，且 $g(w)=f(w^{-1})w^k$。因此
  \[
  H^0(\CP^1,\mathcal{L})=\Span_{\C}\{1,z,\cdots,z^k\},
  \]
  其维数为 $k+1$。这正是 $\mathrm{su}(2)$ 的自旋 $k/2$ 表示。
  我们也可以写出辛形式的具体表达式。在全纯图卡 $U$ 上，
  \[
  \omega=ik\hbar\frac{dz\wedge d\ol z}{(1+\abs{z}^2)^2}.
  \]
\end{example}

\begin{example}{5.4.11}
  我们延续上一个例子，但换一个视角。设 $(\mathcal{L}_0,\nabla_0)$ 是复流形 $M$ 上固定的厄米全纯线丛，其第一陈类给出辛形式
  \[
  \omega=\frac{i}{2\pi}F_{\nabla_0}.
  \]
  考虑第 $k$ 次张量积 $\mathcal{L}=\mathcal{L}_0^{\otimes k}$，带由 $\nabla_0$ 诱导的联络 $\nabla$。则
  \[
  F_\nabla=kF_{\nabla_0}=-i(2\pi k)\omega.
  \]
  于是，若认同
  \[
  \hbar=\frac{1}{2\pi k},
  \]
  则 $(\mathcal{L},\nabla)$ 给出预量子化线丛。
  特别地，半经典极限 $\hbar\to0$ 对应于 $k\to\infty$。
  设 $M$ 紧，于是量子希尔伯特空间 $H^0(M,\mathcal{L})$ 有限维。令
  \[
  \chi(M,\mathcal{L}):=\sum_i(-1)^i\dim H^i(M,\mathcal{L})
  \]
  % [原文 p.292]
  为 $\mathcal{L}$ 的欧拉示性数。由希尔策布鲁赫–黎曼–罗赫定理（Hirzebruch–Riemann–Roch），
  \[
  \chi(M,\mathcal{L})=\int_M\mathrm{ch}(\mathcal{L})\mathrm{td}(M),
  \]
  其中 $\mathrm{td}(M)$ 是 $M$ 的切丛的托德类。在上面的情形中，$\mathcal{L}_0$ 是正线丛，我们有
  \[
  H^i(M,\mathcal{L}_0^{\otimes k})=0
  \qquad\text{对 } i>0 \text{ 且 } k \text{ 充分大}.
  \]
  于是在 $k\to\infty$ 极限下，对 $\mathcal{L}=\mathcal{L}_0^{\otimes k}$ 有首项行为
  \[
  \begin{aligned}
  \dim H^0(M,\mathcal{L})
  &=\int_M e^{\frac{i}{2\pi}F_\nabla}\mathrm{td}(M)
  =\int_M e^{k\omega}\mathrm{td}(M)\\
  &\sim k^n\int_M\frac{\omega^n}{n!}=k^n\vol_\omega(M),
  \end{aligned}
  \]
  其中 $2n=\dim_{\R}M$，$\vol_\omega(M)$ 是 $M$ 的刘维尔体积。在认同
  $\hbar=\frac{1}{2\pi k}$ 之下，我们得到半经典行为
  \[
  \dim H^0(M,\mathcal{L})\sim\frac{1}{(2\pi\hbar)^n}\vol_\omega(M).
  \]
  这引出一个半经典直觉：量子态的数目约等于相空间中以 $\hbar$ 为单位量得的格子数。
\end{example}

\subsection{量子算子}

设 $(M,\omega)$ 是可量子化的辛流形，已选定预量子化线丛 $(\mathcal{L},\nabla)$ 与极化
$P\subset T^\C M$。下面考虑如何把函数 $f\in C^\infty(M)$ 量子化为量子希尔伯特空间
$L^2(M,\mathcal{L})_P$ 上的算子。
有一类特殊的函数可以自然地量子化，即保持极化 $P$ 的那些函数。

\begin{definition}{5.4.12}
  若函数 $f$ 的哈密顿向量场 $V_f$ 满足
  \[
  [V_f,P]\subset P,
  \]
  则称 $f$ 是 $P$-保持的（$P$-preserving）。也就是说，对 $P$ 中任意向量场 $X$，李括号
  $[V_f,X]$ 仍是 $P$ 中的向量场。$P$-保持函数的全体记为 $C^\infty(M)_P$。
\end{definition}

\begin{proposition}{5.4.13}
  设 $\psi\in\Gamma(M,\mathcal{L})$ 是 $\mathcal{L}$ 的 $P$-极化截面，$f\in C^\infty(M)_P$
  是 $P$-保持函数。则 $Q(f)\psi$ 也是 $\mathcal{L}$ 的 $P$-极化截面。
\end{proposition}

\begin{proof}
  回忆预量子化映射
  \[
  Q(f)=i\hbar\nabla_f+f.
  \]

% [原文 p.293]
  设 $X$ 是 $P$ 中任一向量场。则
  \[
  \begin{aligned}
  [\nabla_X,Q(f)]
  &=[\nabla_X,i\hbar\nabla_f+f]\\
  &=i\hbar\nabla_{[X,V_f]}+i\hbar F_\nabla(X,V_f)+X(f)\\
  &=i\hbar\nabla_{[X,V_f]}+\omega(X,V_f)+X(f)\\
  &=i\hbar\nabla_{[X,V_f]}-\iota_X\iota_{V_f}\omega+X(f)\\
  &=i\hbar\nabla_{[X,V_f]}-\iota_Xdf+X(f)\\
  &=i\hbar\nabla_{[X,V_f]}.
  \end{aligned}
  \]
  由假设，$[X,V_f]$ 仍是 $P$ 中的向量场。于是对 $P$-极化截面 $\psi$，
  \[
  \begin{aligned}
  \nabla_XQ(f)\psi
  &=[\nabla_X,Q(f)]\psi
  &&\text{因为 } \nabla_X\psi=0\\
  &=i\hbar\nabla_{[X,V_f]}\psi\\
  &=0
  &&\text{因为 } [X,V_f] \text{ 落在 } P \text{ 中}.
  \end{aligned}
  \]
  因此 $Q(f)\psi$ 也是 $P$-极化截面。
\end{proof}

于是对 $P$-保持函数 $f$，我们可以直接把它量子化为量子希尔伯特空间上的算子 $Q(f)$。
对不是 $P$-保持的函数，量子化要棘手一些。基本想法是：它们的哈密顿流会改变极化，从而把 $P$-极化态变成相对于另一极化的态。然后可以构造联系不同极化对应的量子希尔伯特空间的映射。这就是布拉特纳–科斯坦特–斯滕伯格（Blattner-Kostant-Sternberg，BKS）核构造的想法。我们这里不再深入细节。

\begin{example}{5.4.14}
  考虑相空间
  \[
  \R^{2n},\qquad \omega=\sum_idp_i\wedge dx^i,
  \]
  配以预量子化线丛
  \[
  \nabla=d-\frac{i}{\hbar}\theta,\qquad \theta=\sum_ip_idx^i,
  \]
  与极化
  \[
  P=\Span\left\{\frac{\pa}{\pa p_i}\right\}.
  \]
  考虑函数 $f(\mathbf x,\mathbf p)$，其哈密顿向量场为
  \[
  V_f=\sum_i\left(\pa_{x^i}f\frac{\pa}{\pa p_i}
  -\pa_{p_i}f\frac{\pa}{\pa x^i}\right).
  \]
  由
  \[
  \left[\frac{\pa}{\pa p_k},V_f\right]
  =\sum_i\left(\pa_{p_k}\pa_{x^i}f\right)\frac{\pa}{\pa p_i}
  -\sum_i\left(\pa_{p_k}\pa_{p_i}f\right)\frac{\pa}{\pa x^i}
  \]

% [原文 p.294]
  由此可知，$V_f$ 保持 $P$ 当且仅当
  \[
  \pa_{p_k}\pa_{p_i}f=0,\qquad\forall i,k.
  \]
  于是 $C^\infty(\R^{2n})_P=\left\{a(\mathbf x)+\sum_ib_i(\mathbf x)p_i\right\}$
  由至多关于 $p_i$ 线性的函数组成。这样的函数
  $f=a(\mathbf x)+\sum_ib_i(\mathbf x)p_i$ 量子化为
  \[
  Q(f)=i\hbar\left(V_f-\frac{i}{\hbar}\theta(V_f)\right)+f
  =-i\hbar\sum_ib_i(\mathbf x)\frac{\pa}{\pa x^i}+a(\mathbf x),
  \]
  作为作用在波函数 $\psi=\psi(\mathbf x)$ 上的算子。
\end{example}

\begin{example}{5.4.15}
  考虑相空间
  \[
  \C^n,\qquad \omega=i\pa\ol\pa K,
  \]
  其中 $K$ 是凯勒势。我们取预量子化线丛 $\mathcal{L}$ 为平凡线丛，其厄米内积为
  \[
  \langle f_1,f_2\rangle=\ol f_1f_2e^{-\frac{K}{\hbar}},
  \]
  相容联络为
  \[
  \nabla=d-\frac{1}{\hbar}\pa K.
  \]
  其曲率
  \[
  F_\nabla=-\frac{1}{\hbar}d(\pa K)=-\frac{1}{\hbar}\ol\pa\pa K
  =\frac{1}{\hbar}\pa\ol\pa K=-\frac{i}{\hbar}\omega
  \]
  满足预量子化条件。我们取凯勒极化
  \[
  P=\Span\left\{\frac{\pa}{\pa\ol z^j}\right\},
  \]
  于是 $\mathcal{L}$ 的极化截面就是全纯函数。
  把凯勒形式写成
  \[
  \omega=\sum_{j,m}g_{j\ol m}dz^j\wedge d\ol z^m,
  \]
  其中 $g_{j\ol m}=i\pa_j\ol\pa_{\ol m}K$。记 $g^{\ol mj}$ 为 $g_{j\ol m}$ 的逆矩阵，即
  \[
  \sum_{\ol m}g_{i\ol m}g^{\ol mj}=\delta_i{}^j.
  \]
  设 $f$ 是 $\C^n$ 上的函数。其哈密顿向量场为
  \[
  V_f=\sum_{j,\ol m}\left(g^{\ol mj}\pa_{\ol m}f\frac{\pa}{\pa z^j}
  -g^{\ol mj}\pa_jf\frac{\pa}{\pa\ol z^m}\right),
  \]
  其中 $\pa_{\ol m}=\frac{\pa}{\pa\ol z^m}$，$\pa_j=\frac{\pa}{\pa z^j}$。由于
  \[
  \left[\frac{\pa}{\pa\ol z^i},V_f\right]
  =\sum_{j,\ol m}\frac{\pa}{\pa\ol z^i}(g^{\ol mj}\pa_{\ol m}f)\frac{\pa}{\pa z^j}
  \qquad\text{模 } P,
  \]

% [原文 p.295]
  由此可知，$f$ 保持 $P$ 当且仅当
  \[
  \sum_{\ol m}\frac{\pa}{\pa\ol z^i}
  \left(g^{\ol mj}\frac{\pa}{\pa\ol z^m}f\right)=0,\qquad\forall i,j.
  \]
  现在取特殊情形
  \[
  K=|\mathbf z|^2,\qquad g_{j\ol m}=i\delta_{jm}.
  \]
  这时 $P$-保持条件变成
  \[
  \frac{\pa}{\pa\ol z^i}\frac{\pa}{\pa\ol z^j}f=0,\qquad\forall i,j,
  \]
  即 $f$ 至多关于 $\ol z^i$ 线性。若再要求 $f$ 取实值，则仅有的实 $P$-保持函数形如
  \[
  f=a+b_kz^k+\ol b_k\ol z^k+c_{k\ol j}z^k\ol z^j,
  \]
  其中 $a\in\R$，$b_k,c_{k\ol j}\in\C$，且 $\ol{c_{k\ol j}}=c_{j\ol k}$。$f$ 的哈密顿向量场为
  \[
  V_f=-i\sum_k\left(\pa_kf\frac{\pa}{\pa\ol z^k}
  -\ol{\pa_kf}\frac{\pa}{\pa z^k}\right).
  \]
  于是 $f$ 量子化为
  \[
  \begin{aligned}
  Q(f)&=i\hbar\left(V_f-\frac{1}{\hbar}\pa K(V_f)\right)+f\\
  &=\hbar\sum_k\pa_{\ol k}f\frac{\pa}{\pa z^k}
  -\sum_k\ol z^k\frac{\pa}{\pa\ol z^k}f+f\\
  &=a+\sum_kb_kz^k+\hbar\sum_k\left(\ol b_k+c_{j\ol k}z^j\right)
  \frac{\pa}{\pa z^k},
  \end{aligned}
  \]
  作为作用在全纯函数上的算子。

  这使得全纯表示对二次哈密顿量的量子化格外好用。例如考虑一维谐振子，哈密顿量为
  \[
  \mathscr{H}=|z|^2\qquad\text{于 } \C.
  \]
  上述凯勒极化把 $\mathscr{H}$ 量子化为算子
  \[
  Q(\mathscr{H})=\hbar z\frac{\pa}{\pa z},
  \]
  它作用在 Segal–Bargmann 空间上。注意
  \[
  Q(\mathscr{H})z^n=n\hbar z^n,
  \]
  因此 $Q(\mathscr{H})$ 的本征值是以 $\hbar$ 为单位的非负整数。然而这并不是正确的谱：马斯洛夫修正缺失了。
  这与我们在半经典量子化的讨论中遇到的问题相同。出路是修改量子希尔伯特空间，把半密度也包含进来。

% [原文 p.296]
  在凯勒量子化中，这一修改相当于把量子态 $\psi$ 定义为
  $\mathcal{L}\otimes\sqrt{K_M}$ 的截面，其中 $K_M$ 是典范线束，$\sqrt{K_M}$ 是它的平方根。
  在上述谐振子的例子里，这一修改相当于把态
  \[
  z^n\longrightarrow z^n\sqrt{dz},
  \]
  此后它在 $Q(\mathscr{H})=\hbar z\frac{\pa}{\pa z}$ 之下就有正确的本征值
  $(n+\frac12)\hbar$。
\end{example}

\takeaway{本节由狄拉克量子化原理指出经典可观测量的映射无法满足全部要求，格罗嫩沃尔德–范霍夫定理给出精确形式；预量子化线丛与极化把哈密顿向量场提升为截面算子，陈类的整性条件刻画可量子化的辛流形。}

\chapterreview{本章讨论几何量子化。§5.1 整理相空间的辛几何：辛形式、哈密顿向量场、泊松括号与拉格朗日子流形。§5.2 讲半经典量子化：玻尔–索末菲条件、刘维尔测度与马斯洛夫修正，给出量子态与相空间体积的对应。§5.3 由一维谐振子引出海森堡群，构造其李代数与群结构，在其上建立薛定谔表示与巴格曼–福克表示，借斯通–冯·诺依曼定理说明表示的唯一性，并给出 Segal–Bargmann 变换。§5.4 讲几何量子化：狄拉克量子化原理及其不可能性、预量子化线丛与可量子化条件、极化的定义与量子希尔伯特空间、量子算子与半密度修正。}

